\documentclass[11pt,a4paper]{article}
\usepackage[utf8]{inputenc}
\usepackage[T1]{fontenc}
\usepackage[italian]{babel}
\usepackage{amsmath}
\usepackage{amssymb}
\usepackage[margin=1.8cm]{geometry}
\usepackage{xcolor}
\usepackage{booktabs}
\usepackage{hyperref}
\usepackage{tcolorbox}

\definecolor{primary}{RGB}{30, 60, 110}
\definecolor{accent}{RGB}{200, 70, 40}
\definecolor{success}{RGB}{35, 125, 60}

\title{\textbf{\Large Specifiche di Stampa 3D — Ricambi Korg SP-200}\\\large Leggio Modulare v4 (600\,mm) e Copri-Fader parametrico}
\author{\textbf{Istruzioni per Creality Ender 3 V2 e Montaggio}}
\date{\today}

\begin{document}
\maketitle

\begin{tcolorbox}[colback=primary!5!white,colframe=primary,title=\textbf{Sintesi del Progetto — Versione v4 per Creality Ender 3 V2}]
I modelli sono calibrati sulle misure reali del pianoforte digitale \textbf{Korg SP-200} e dimensionati specificamente per il piatto di stampa della \textbf{Creality Ender 3 V2} ($220\times 220$\,mm con piano in vetro carborundum).

\textbf{Caratteristiche e Migliorie della Versione v4}:
\begin{itemize}
    \item \textbf{Larghezza da concerto estesa a 600\,mm} ($+15$\,cm rispetto alla v1): il leggio sporge di $85$\,mm per lato oltre i perni di aggancio (interasse fisso a $430$\,mm), sostenendo comodamente libri aperti e spartiti a tre facciate.
    \item \textbf{Architettura a 3 pezzi da 200\,mm}: corpo centrale continuo (con la rientranza a V solida e senza giunture visibili al centro) e due ali laterali. Ciascun pezzo misura $200 \times 187.5$\,mm, lasciando ampio margine dai bordi e dalle mollette del piatto in vetro.
    \item \textbf{Inclinazione anti-spanciamento a $24^\circ$}: sfrutta la gravità per mantenere i fogli A4 singoli ben adesi allo schienale.
    \item \textbf{Mensolina profonda 33\,mm}: ampio alloggiamento per album rilegati e spessi.
    \item \textbf{Parapetto svasato a $105^\circ$}: ribassato a $10$\,mm con cresta arrotondata per voltare le pagine senza impuntamenti.
\end{itemize}
\end{tcolorbox}

\section{Materiale e Parametri di Stampa Ender 3 V2}
\begin{itemize}
    \item \textbf{Filamento}: \textbf{PLA o PETG} (consigliato Nero opaco / Black Matte o Grafite scuro).
    \item \textbf{Piatto in vetro Ender 3 V2}: Temperatura piatto a \textbf{$60^\circ$C} per PLA ($70^\circ$C per PETG). Pulire bene il vetro con alcool isopropilico o acqua calda e sapone per piatti prima della stampa per garantire perfetta aderenza.
    \item \textbf{Temperatura Ugello}: $200\text{--}205^\circ$C (PLA) / $230\text{--}235^\circ$C (PETG).
\end{itemize}

\section{Elenco File STL da Stampare}
\begin{table}[h!]
\centering
\begin{tabular}{@{}llcc@{}}
\toprule
\textbf{File STL} & \textbf{Descrizione} & \textbf{Q.tà} & \textbf{Tempo} \\
\midrule
\texttt{leggio\_provino\_test.stl} & \textbf{Test rapido d'inserimento nel piano} & \textbf{1 pezzo} & \textbf{$\sim$12 min} \\
\texttt{copri\_fader.stl} & Copri-fader volume/brilliance & 2 pezzi & $\sim$12 min \\
\texttt{leggio\_centro.stl} & Corpo centrale da 200\,mm (V continua) & 1 pezzo & $\sim$1h 45m \\
\texttt{leggio\_ala\_sinistra.stl} & Ala sinistra da 200\,mm (con perno Korg) & 1 pezzo & $\sim$1h 45m \\
\texttt{leggio\_ala\_destra.stl} & Ala destra da 200\,mm (con perno Korg) & 1 pezzo & $\sim$1h 45m \\
\texttt{leggio\_spine\_giunzione.stl} & Set di 6 spine cilindriche da 22\,mm & 1 set (6 spine) & $\sim$5 min \\
\bottomrule
\end{tabular}
\end{table}

\begin{tcolorbox}[colback=success!5!white,colframe=success,title=\textbf{Stampa di Prova Rapida: Provino di Test (12 Minuti)}]
Prima di stampare i pezzi grandi, \textbf{è sempre consigliato stampare il file \texttt{leggio\_provino\_test.stl}}:
\begin{itemize}
    \item \textbf{Cosa contiene}: Il perno orizzontale a L da 26\,mm, il montante di battuta e i primi 9\,cm di mensolina da 33\,mm con parapetto svasato a $105^\circ$. Consuma meno di 5 grammi di filo.
    \item \textbf{Verifica sul pianoforte}: Inserire il perno nel foro posteriore del pianoforte: deve scorrere dritto per tutti i 26\,mm arrestandosi a battuta con la boccola Korg, con la mensolina radente a circa $1.5\text{--}2$\,mm dalla carrozzeria.
\end{itemize}
\end{tcolorbox}

\section{Impostazioni Slicer Consigliate (Cura / PrusaSlicer / Creality Print)}

\subsection{1. Copri-Fader (\texttt{copri\_fader.stl})}
\begin{itemize}
    \item \textbf{Layer}: $0.12$\,mm o $0.16$\,mm. Base piatta sul piatto ($Z=0$), tacca verso l'alto. Supporti \textbf{disattivati}.
    \item \textbf{Pareti}: 4 perimetri solidi.
\end{itemize}

\subsection{2. Tre Parti Principali del Leggio (Centro e Ali Laterali)}
\begin{itemize}
    \item \textbf{Altezza Layer}: $0.20$\,mm.
    \item \textbf{Pareti (Wall loops / Perimeters)}: \textbf{4 o 5 pareti} (\textit{fondamentale per la rigidità strutturale dei tubolari da 9\,mm}).
    \item \textbf{Riempimento (Infill)}: $20\%$--$25\%$ con motivo \textbf{Giroide (Gyroid)}.
    \item \textbf{Orientamento sul piatto}: Sdraiato piatto sul dorso/schiena (usare la funzione \textit{"Appoggia sulla faccia"} / \textit{"Place on face"}).
    \item \textbf{Posizionamento su Ender 3 V2}: Centrare ciascun pezzo sul piatto; le dimensioni di $200 \times 187.5$\,mm lasciano $10$\,mm di spazio sui lati, evitando le mollette anteriori e posteriori del vetro.
    \item \textbf{Supporti}: Supporti ad albero (\textit{Tree supports}) raccomandati solo per il perno a gomito posteriore sulle due ali.
\end{itemize}

\subsection{3. Spine di Giunzione (\texttt{leggio\_spine\_giunzione.stl})}
\begin{itemize}
    \item \textbf{Geometria}: Il file contiene 6 spine cilindriche identiche da $22$\,mm ($\varnothing\,4.6$\,mm, con smusso a $45^\circ$).
    \item \textbf{Tolleranze}: I fori femmina sul leggio sono alesati a $\varnothing\,5.4$\,mm e profondi $18$\,mm per lato: le spine scorrono senza attriti e lasciano $7$\,mm di aria cieca a fondo foro (impossibile che vadano a pacco).
    \item \textbf{Riempimento}: $100\%$ o 4 pareti (funzione di tiranti strutturali). Sdraiate orizzontalmente sul piatto.
\end{itemize}

\section{Istruzioni di Assemblaggio e Montaggio}
\begin{enumerate}
    \item Rimuovere eventuali residui di supporto dai perni posteriori delle due ali.
    \item \textbf{Prova a secco}: inserire 3 spine nelle sedi di sinistra del corpo centrale e unire l'ala sinistra; inserire 3 spine nelle sedi di destra e unire l'ala destra. Verificare che i tre pezzi chiudano a filo perfetto.
    \item Applicare alcune gocce di colla cianoacrilica (tipo Super Attak) nei fori e sulle superfici di contatto delle giunzioni.
    \item Unire e premere saldamente i componenti per 30--60 secondi.
    \item (\textit{Opzionale}): Applicare due piccoli feltrini adesivi da 1\,mm sotto la mensolina per un contatto morbido con la carrozzeria del pianoforte.
\end{enumerate}

\end{document}
